\begin{exercise}{36}
The goal of this exercise is to show that any essentially surjective and fully
faithful functor $F\colon\mathcal C\to\mathcal D$ between ordinary categories is an equivalence.
\begin{enumerate}[label=(\arabic*),leftmargin=*,itemsep=0pt,parsep=0pt,
                  topsep=.5\baselineskip]
 \item Show that $F$ admits an adjoint $G$. Hint: Use Exercise~\ref{ex:10}.
 \item Show that $G$ is itself fully faithful.
 \item Show that an adjoint pair $(F,G)$ of fully faithful functors makes $F$ an equivalence with
       $G$ an inverse.
\end{enumerate}

\begin{proof}
For an adjunction $F\dashv G$, write $\eta$ for the unit and
$\varepsilon$ for the counit. Lemma~1.2.19 says that $F$ is
fully faithful exactly when $\eta$ is an isomorphism, and $G$
is fully faithful exactly when $\varepsilon$ is an isomorphism.

(1) By essential surjectivity, choose for each $d\in\mathcal D$
an object $Gd\in\mathcal C$ and an isomorphism
$\varepsilon_d\colon FGd\to d$.
The maps
\[
 \operatorname{Hom}_{\mathcal C}(c,Gd)
 \xrightarrow{F}\operatorname{Hom}_{\mathcal D}(Fc,FGd)
 \xrightarrow{\varepsilon_d\circ-}
 \operatorname{Hom}_{\mathcal D}(Fc,d)
\]
are bijective: the first arrow uses full faithfulness, and
the second postcomposes with an isomorphism.
Exercise~\ref{ex:10} therefore supplies a right adjoint $G$
whose counit has the chosen components $\varepsilon_d$.

(2) This counit is an isomorphism, so $G$ is fully faithful
by the criterion above.

(3) For any adjoint pair of fully faithful functors, the same
criterion makes both $\eta$ and $\varepsilon$ natural
isomorphisms. They exhibit $G$ as a quasi-inverse of $F$,
so $F$ is an equivalence.
\end{proof}
\end{exercise}
